\documentclass[11pt, a4paper]{article}

% ── Packages ──────────────────────────────────────────────────────────────────
\usepackage[T1]{fontenc}
\usepackage[utf8]{inputenc}
\usepackage{lmodern}
\usepackage{microtype}
\usepackage{amsmath, amssymb, amsthm}
\usepackage{mathtools}
\usepackage{nicefrac}
\usepackage{geometry}
\usepackage{hyperref}
\usepackage{cleveref}
\usepackage{enumitem}
\usepackage{booktabs}
\usepackage{listings}
\usepackage{xcolor}
\usepackage[most]{tcolorbox}
\tcbuselibrary{listings, breakable}
\usepackage{biblatex}
\addbibresource{refs.bib}
\usepackage{xspace}

\geometry{margin=2.5cm}

% ── Theorem environments ───────────────────────────────────────────────────────
\theoremstyle{plain}
\newtheorem{theorem}{Theorem}[section]
\newtheorem{proposition}[theorem]{Proposition}
\newtheorem{lemma}[theorem]{Lemma}
\newtheorem{corollary}[theorem]{Corollary}

\theoremstyle{definition}
\newtheorem{definition}[theorem]{Definition}
\newtheorem{example}[theorem]{Example}
\newtheorem{remark}[theorem]{Remark}

% ── Notation shorthands ───────────────────────────────────────────────────────
\newcommand{\catena}{\textsf{catena}\xspace}

\newcommand{\Z}{\mathbb{Z}}
\newcommand{\Q}{\mathbb{Q}}
\newcommand{\R}{\mathbb{R}}
\newcommand{\N}{\mathbb{N}}

\newcommand{\Zp}{\mathbb{Z}^+}  % positive integers
\newcommand{\Qp}{\mathbb{Q}^+}  % positive rationals
\newcommand{\Rp}{\mathbb{R}^+}  % positive reals

\newcommand{\scf}[1]{[#1]}                    % [a0; a1, a2, ...]  -- inline SCF
\renewcommand{\cfrac}[2]{\dfrac{#1}{#2}}      % override amsmath \cfrac for display-size fractions
\newcommand{\convergent}[1]{x_{#1}}           % n-th convergent
\newcommand{\pq}[1]{(p_{#1}, q_{#1})}         % convergent numerator/denominator pair
\newcommand{\abs}[1]{\left\lvert #1 \right\rvert}

% ── Code listings style ────────────────────────────────────────────────────────
\lstdefinestyle{python}{
  language=Python,
  basicstyle=\ttfamily\small,
  keywordstyle=\color{blue!70!black}\bfseries,
  commentstyle=\color{gray},
  stringstyle=\color{green!50!black},
  numbers=left, numberstyle=\tiny\color{gray},
  frame=none, breaklines=true,
  aboveskip=0pt, belowskip=0pt,
}

% ── catenadef environment ──────────────────────────────────────────────────────
% Usage:
%   \begin{catenadef}[optional title]
%     \begin{lstlisting}[style=python]
%       ...
%     \end{lstlisting}
%   \end{catenadef}
%
% Visually pairs with the immediately preceding definition/theorem block.
% The optional argument overrides the default header label.
\newtcolorbox{catenadef}[1][]{%
  enhanced, breakable,
  colback=gray!6,
  colframe=gray!30,
  leftrule=3pt,
  rightrule=0.4pt,
  toprule=0.4pt,
  bottomrule=0.4pt,
  arc=1.5pt,
  top=4pt, bottom=4pt, left=6pt, right=6pt,
  before upper={%
    \textcolor{gray!70!black}{%
      \textsf{\small\textbf{\catena}}%
      \ifx\relax#1\relax\else{\textcolor{gray!60!black}{\small~—~#1}}\fi%
    }\par\smallskip%
  },
  before skip=2pt,
  after skip=6pt,
}

% ── Metadata ──────────────────────────────────────────────────────────────────
\title{\textbf{catena} \\[0.5em]
       \large Architecture, Algorithm Decisions, and Formal Theory \\[0.5em]
       \large \textit{v0.1, Draft}}
\author{Lorenzo Silva Moore\footnote{Facultad de Ciencias (Faculty of Sciences), UNAM, México.\\\indent\indent \texttt{lorenzosilvamoore@ciencias.unam.mx} $\odot$ \texttt{lorenzo.smb.rayo@gmail.com}}}
\date{\today}

% ══════════════════════════════════════════════════════════════════════════════
\begin{document}
\maketitle
\tableofcontents
\bigskip
\hrule
\bigskip

% ─────────────────────────────────────────────────────────────────────────────
\section{Introduction and Scope}
% ─────────────────────────────────────────────────────────────────────────────

This document records the mathematical foundations and algorithmic design decisions
behind \catena, a Python library for simple continued fractions.
It is intended as a living document: new sections will be added as the library
grows to cover further aspects of the theory.

The current implementation (\texttt{v0.4.0}) covers finite SCFs over the rationals,
infinite generative SCFs, and periodic SCFs for quadratic irrationals.

\subsection{What continued fractions offer}

A real number can be represented in several ways: as a floating-point approximation,
as an exact rational via \texttt{fractions.Fraction}, as a fixed-precision decimal,
or as a continued fraction.  Each representation makes different trade-offs.

Floats are fast and ubiquitous, but they introduce rounding error at every operation.
A float is not a rational number; it is the nearest representable binary fraction,
and the gap between the two is invisible unless you care to look for it.

\texttt{fractions.Fraction} is exact and supports arbitrary precision, but it
accumulates large numerators and denominators quickly.  Adding two fractions
$p_1/q_1 + p_2/q_2$ yields $(p_1 q_2 + p_2 q_1)/(q_1 q_2)$; in the worst case
(coprime denominators) the denominator grows from $n$ digits to $2n$ digits before
any reduction, and every operation incurs a GCD computation whose cost grows with
digit count.

Continued fractions offer a different kind of exactness.  The partial quotients
$[a_0;\, a_1, a_2, \ldots]$ are small integers — typically single digits for
naturally occurring constants — so the representation is compact.  The $n$-th
convergent $p_n/q_n$ is automatically in lowest terms by the determinant identity
$p_n q_{n-1} - p_{n-1} q_n = \pm 1$, with no GCD step required.  Convergents
also have an optimality property: $p_n/q_n$ is the best rational approximation
to the true value among all fractions with denominator $\leq q_n$.

For irrational numbers, the continued fraction is the natural representation.
$\sqrt{2} = [1;\overline{2}]$, $\varphi = [1;\overline{1}]$, and more generally
every quadratic irrational has an eventually periodic expansion.  Floats and
decimals can only approximate; the SCF records the structure exactly.

These properties make SCFs useful for problems where rational approximation
quality matters — gear-ratio design, modular arithmetic, Pell's equation, and
the analysis of approximation algorithms, among others.  \catena is built around
this perspective: the library works with the SCF itself as the primary object,
deriving convergents and numeric approximations from it on demand.

\subsection{Scope}

\catena provides three main abstractions:

\begin{itemize}[nosep]
  \item \texttt{SimpleContinuedFraction} — an infinite SCF $[a_0;\,a_1,a_2,\ldots]$
        driven by an arbitrary callable generator.  Convergents are memoised as they
        are requested; the cache is shared when the same tail is viewed with a different
        integer part (e.g.\ after an integer shift).
  \item \texttt{FiniteSimpleContinuedFraction} — a finite SCF
        $[a_0;\,a_1,\ldots,a_k]$ over $\Q$, constructible from a rational, float,
        or decimal string.  Supports exact rational arithmetic and multiplicative inverse.
  \item \texttt{PeriodicSimpleContinuedFraction} — an eventually-periodic SCF
        $[a_0;\,a_1,\ldots,a_m,\,\overline{b_1,\ldots,b_r}]$ representing a quadratic
        irrational $(P + \sqrt{D})/Q$.  Provides algebraic operations: conjugate,
        inverse, negation, and exact decimal evaluation.
\end{itemize}

\begin{remark}
  \texttt{PeriodicSimpleContinuedFraction} is listed here for completeness.
  Its full mathematical treatment — Lagrange's theorem on quadratic irrationals,
  the Pell equation, and the algebraic operations
  (\texttt{conjugate}, \texttt{inverse}, \texttt{quadratic\_coefficients}) —
  is deferred to a forthcoming section.
\end{remark}

The mathematical content of this document is therefore structured around these
three cases.  It is not intended to be an exhaustive treatment of continued fraction
theory; for that, see \cite{khinchin1964} and \cite{rockett1992}.  The goal is to
give enough background to justify the algorithms and design choices in the library.

\subsection{Conventions}
      
Throughout this document:
\begin{itemize}[nosep]
      \item $\N$ denotes the natural numbers (including zero).
      \item For any $n \in \N$, we implicitly assume the set-theoretic construction $n = \{0, 1, \ldots, n-1\}$.
      \item $\Z$, $\Q$, $\R$ denote the integers, rationals, and reals respectively.
      \item Likewise, $\Zp$, $\Qp$, $\Rp$ denote the positive (non-zero) integers, rationals, and reals respectively.
      \item $[a_0;\, a_1, a_2, \ldots, a_n]$ denotes a finite SCF with $a_0 \in \Z$ and
            $a_k \in \Zp$ for $k \geq 1$.
      \item $c_k = (p_k, q_k)$ denotes the $k$-th convergent numerator--denominator pair, with
            $q_k \in \Zp$ by convention.
      \item The function $\{\cdot\}: \R \to \R$ denotes the fractional part, so $\{x\} = x - \lfloor x \rfloor$.
\end{itemize}

\begin{remark}
  We follow the set-theoretic convention $\N = \{0, 1, 2, \ldots\}$ including zero.
  The positive integers $\Zp = \{1, 2, 3, \ldots\}$ exclude zero.
  In particular, all body partial quotients $a_k$ (for $k \geq 1$) lie in $\Zp$,
  while the integer part $a_0 \in \Z$ may be zero or negative.
\end{remark}

% ─────────────────────────────────────────────────────────────────────────────
\section{Simple Continued Fractions\texorpdfstring{\protect\footnotemark}{}}\footnotetext{As with all the other sections, the following is not meant to be an exhaustive 
treatment of the topic, but rather a sketch of the main results and algorithms related to SCFs that are relevant to \catena. 
For a more comprehensive treatment, see \cite{khinchin1964} and \cite{rockett1992}.
}
% ─────────────────────────────────────────────────────────────────────────────

In the most general sense, a continued fraction is an expression of the form
\[
  a_0 + \cfrac{b_1}{a_1 + \cfrac{b_2}{a_2 + \cfrac{b_3}{\ddots + \cfrac{b_n}{a_n} + \ldots}}}
\]
where $a_0, a_1, \ldots, a_n$ and $b_1, b_2, \ldots, b_n$, depending on the context, 
may be assumed to be integers, rational numbers, either real numbers or complex numbers,
and even functions. For the purposes of \catena, and by extension this document, we 
we restrict attention to the so-called \textit{simple continued fractions} (SCFs), that is, those 
where $b_n = 1$ for all $n$ and $a_n$ are integers with $a_1, a_2, \ldots, a_n, \ldots > 0$.\footnote{
    In the literature is common to allow $a_0 \in \R$ (for instance, see \cite{khinchin1964}), 
    but for the sake of avoiding floating-point arithmetic, as well as to gain a notion of uniqueness,
    we will restrict it to $\Z$. If nothing more, we do not lose much generality since any real number 
    admits a unique SCF expansion with integer part $a_0 = \lfloor x \rfloor$ and fractional part 
    $x - a_0 = 1/(a_1 + 1/(a_2 + \ldots))$. While we won't stop to prove it here, 
    it is worth noting that the SCF expansion of a rational number is finite, while that of an 
    irrational number is infinite.
} Namely, we are left with expressions of the form
\[
  a_0 + \cfrac{1}{a_1 + \cfrac{1}{a_2 + \cfrac{1}{\ddots + \cfrac{1}{a_n} +\cdots}}}
\]
with $a_0 \in \Z$ and $a_1, a_2, \ldots, a_n, \ldots \in \Zp$. For brevity, it is common to write this 
expression as $[a_0;\, a_1, a_2, \ldots, a_n, \ldots]$. The $a_n$ are called the \textit{partial quotients} 
of the SCF, and the $n$-th \textit{convergent} is the rational number obtained by truncating the 
expansion at $a_n$, that is, $[a_0;\, a_1, a_2, \ldots, a_n]$.

As any sequence of indexable objects, one could represent an SCF as a function $f: \N \to \Z$ such that 
$f(0) \in \Z$ and $f(n) \in \Zp$ for $n > 0$, thereby encoding the entire expansion uniformly within a 
single map. However, as suggested by the standard notation $[a_0;\, a_1, a_2, \ldots]$, it is often 
conceptually and computationally advantageous to distinguish and separate the integer part from the 
remaining partial quotients.

Accordingly, in this work we adopt a split representation: we treat the first partial quotient $a_0 \in \Z$ 
as the \textit{integer part} of the SCF, and encode the remaining coefficients via a generator $f: \N \to \Zp$ 
where $f(n) = a_{n+1}$. This distinction is not merely notational; it reflects how SCFs are 
constructed and manipulated algorithmically in \catena, and will play a central role in the 
implementations developed in the following sections.

\subsection{Core concepts and implementations}

\begin{definition}[Generating Function]
  Let $I \in \N \cup \{N\}$. A \textit{generator} indexed by $I$ is any function 
  $g: I \to \Zp$. 
  
  When $I = \N$, we say that $g$ is an \textit{infinite generator}, and when 
  $I \in \N$, we say that $g$ is a \textit{finite generator}. Furthermore, in the finite case, we can say that $g$ is 
  a I\textit{th-order} generator.
\end{definition}

\begin{remark}
    Note we allow $I = 0$. In this case, the generator is vacuous and the corresponding SCF reduces to $[a_0]$ for some $a_0 \in \Z$.
\end{remark}

\begin{catenadef}[\texttt{Generator}, \texttt{CachedGenerator}, \texttt{FiniteGenerator}]
\begin{lstlisting}[style=python]
from catena.generators import Generator, CachedGenerator, FiniteGenerator

# Infinite generator: a_n = 2 for all n > 0, with integer part a_0 = 1 (\sqrt{2} = [1; 2, 2, 2, ...])
gen = Generator(lambda n: 2)

# CachedGenerator is useful when f(n) is expensive to evaluate.
# Here, naive recursive Fibonacci is O(2^n) per call without memoisation;
# wrapping it in CachedGenerator ensures each index is computed at most once
# across all calls.
def fib(n):
    return fib(n - 1) + fib(n - 2) if n > 1 else n + 1

cgen = CachedGenerator(fib)   # fib(n) is called once per n, then cached
cgen(10) # computes fib(10) and caches results for fib(0) through fib(10)
cgen(15) # computes fib(11) through fib(15), reusing cached results as needed

# Finite generator: fixed sequence [7, 15, 1, 292, ...]  (pi = [3; 7, 15, 1, 292, ...])
fgen = FiniteGenerator([7, 15, 1, 292, 1, 1, 1, 2]) # the well-known SCF expansion of \pi begins thus
fgen(0) # 7
fgen(3) # 292
fgen(10) # raises IndexError

\end{lstlisting}
\end{catenadef}


\begin{definition}[Simple Continued Fraction]
  A \textit{simple continued fraction} is an expression of the form
  \[
    x = a_0 + \cfrac{1}{a_1 + \cfrac{1}{a_2 + \cfrac{1}{\ddots + \cfrac{1}{a_n} + \cdots}}}
  \]
  where $a_0$ is an integer and $a_1, a_2, \ldots, a_n, \ldots$ are positive integers. Thus,
  given $a_0 \in \Z$ and a generator $g: I \to \Zp$, such that $g(n) = a_{n+1}$, we represent the SCF as
  the ordered pair $x = (a_0, g)$.
\end{definition}

\begin{catenadef}[\texttt{SimpleContinuedFraction}, \texttt{FiniteSimpleContinuedFraction}]
\begin{lstlisting}[style=python]
from catena import SimpleContinuedFraction, FiniteSimpleContinuedFraction

# Infinite SCF: golden ratio phi = [1; 1, 1, 1, ...]
phi = SimpleContinuedFraction(lambda n: 1, integer_part=1)

# Finite SCF: 355/113 = [3; 7, 16]
scf = FiniteSimpleContinuedFraction.from_rational((355, 113))
scf.integer_part # 3
scf.generator    # FiniteGenerator([7, 16])
\end{lstlisting}
\end{catenadef}

\begin{definition}[Segment]
    Let $x = (a_0, g)$ be a simple continued fraction, where $g: I \to \Zp$ is a generator. We 
    define the \textit{segment} of $x$ at index $n \in I$ as the SCF obtained by truncating the first 
    $n$ partial quotients from $g$. 
    
    Formally, for $n \in I$, the segment of $x$ at index $n$ is
    given by $x_n = (a_0, g_n)$, where $g_n:=g|_n$ is the restriction of $g$ to $n$.
\end{definition}

\begin{catenadef}[\texttt{SimpleContinuedFraction.segment}]
\begin{lstlisting}[style=python]
from catena import SimpleContinuedFraction

# Example: for phi = [1; 1, 1, 1, ...], the segment at index n is also phi
phi = SimpleContinuedFraction(lambda n: 1, integer_part=1)
phi.segment(0) # [1;]
phi.segment(1) # [1; 1]
phi.segment(10) # [1; 1, 1, 1, 1, 1, 1, 1, 1, 1, 1]
\end{lstlisting}
\end{catenadef}

\begin{remark}
    Given that the segment is defined for $n \in I$, for finite SCFs, the segment is only defined 
    for $n \leq \text{size}$. For $n > \text{size}$, the segment is not defined since it would 
    require access to partial quotients that do not exist in the finite generator. Therefore, 
    in the finite case, calling \texttt{segment(n)} with $n > \text{size}$ will raise an \texttt{IndexError}.

    As for \texttt{segment(size)}, \texttt{self.shift(0)} is returned since the segment at index \texttt{size} 
    corresponds to the original SCF itself, as to avoid unnecessary copying.
\end{remark}

\begin{definition}[Remainder]
    Let $x = (a_0, g)$ be a simple continued fraction, where $g: I \to \Zp$. 
    For $n \in I$, we define the \textit{remainder} of $x$ at index $n$ as the SCF
    \[
        r_n(x) = (g(n), h),
    \]
    where $h: |I \setminus (n+1)| \to \Zp$ is defined by
    \[
        h(m) = g(m + (n+1)).
    \]
\end{definition}

\begin{remark}
    As of now, \catena does not have a dedicated implementation for remainders. We mention them 
    here because they will come up in some of the theoretical results.
\end{remark}

\begin{definition}[Convergent]
    Let $x = (a_0, g)$ be a simple continued fraction, where $g: I \to \Zp$. 
    For $n \in I$, we define the $n$-th \textit{convergent} of $x$ as the rational number
    \[
        c_n = \frac{p_n}{q_n}
    \]
    given by the value of the segment $x_n$.
\end{definition}

\begin{remark}[Segment vs Convergent]
    The segment $x_n$ is itself a (finite) SCF, while the convergent $\nicefrac{p_n}{q_n}$ 
    is a rational number. They are related by evaluation: $x_n$ represents 
    $\nicefrac{p_n}{q_n}$, and every rational number admits a unique finite SCF expansion 
    of this form.
\end{remark}

\begin{definition}[Tail]
    Let $x = (a_0, g)$ be a simple continued fraction, where $g: I \to \Zp$. For $n \in I$, the 
    \textit{tail} of $x$ is the SCF given by $t(x) = (0, g)$. This is to say, the tail of $x$ is the SCF 
    corresponding to the fractional part of $x$, $\{x\}$.

    Likewise, the n\textit{-th tail convergent} of $x$ is the $n$-th convergent of $t(x)$.
\end{definition}

\begin{catenadef}[\texttt{SimpleContinuedFraction.tail}]
\begin{lstlisting}[style=python]
from catena import SimpleContinuedFraction

# Example: for phi = [1; 1, 1, 1, ...], the tail is also phi
phi = SimpleContinuedFraction(lambda n: 1, integer_part=1)
phi.tail() # [0; 1, 1, 1, ...]

phi.tail_convergent(7) # 21/34
phi.tail().convergent(7) # 21/34
\end{lstlisting}
\end{catenadef}

\begin{theorem}[Recurrence Relations for Tail Convergents\footnote{
    See, for instance, \cite{khinchin1964}, Theorem 1, p. 4.
    For a more general version of this result, see \cite{jonesthron1980}, Theorem 2.1, p. 22.
    }]
    Let $x = (a_0, g)$ be a simple continued fraction, where $g: I \to \Zp$. 
    For $n \in I$, define the $n$-th tail convergent by
    \[
        \frac{h_n}{k_n}
        =
        [0;\, g(0), g(1), \ldots, g(n)].
    \]
    Then the numerators and denominators satisfy the recurrence relations
    \[
        h_n = g(n)\, h_{n-1} + h_{n-2}, \quad 
        k_n = g(n)\, k_{n-1} + k_{n-2} \qquad (n \geq 0),
    \]
    with initial conditions
    \[
        h_{-2} = 1,\quad h_{-1} = 0, \qquad
        k_{-2} = 0,\quad k_{-1} = 1.
    \]
\end{theorem}

\begin{remark}
  We extend the indexing of the sequences $(h_n)$ and $(k_n)$ to include 
  the values $n = -2$ and $n = -1$, with
  \[
      h_{-2} = 1,\quad h_{-1} = 0, \qquad
      k_{-2} = 0,\quad k_{-1} = 1.
  \]
  These values should be understood as formal initial conditions for the 
  recurrence, rather than as convergents. This convention allows the 
  recurrence relations to hold uniformly for all $n \geq 0$ without requiring 
  special cases, which is convenient for implementation.
\end{remark}

\begin{corollary}[Convergent from Tail Convergents]
    Let $x = (a_0, g)$ be a simple continued fraction, and let 
    $\frac{h_n}{k_n} = [0;\, g(0), \ldots, g(n)]$ be the $n$-th tail convergent. 
    Then the $n$-th convergent of $x$ is given by
    \[
        c_n = \frac{p_n}{q_n} = a_0 + \frac{h_n}{k_n}
        = \frac{a_0 k_n + h_n}{k_n}.
    \]
\end{corollary}

\begin{remark}[Caching Tail Convergents]
    The recursive nature of the tail convergents makes them ideal candidates for caching. 
    By storing previously computed tail convergents, we can avoid redundant calculations 
    and significantly improve the efficiency of algorithms that rely on them. A pretty direct
    example would be the scenario of iterating over the convergents of a given SCF: 
    by caching the tail convergents, we can compute each new convergent in constant time after 
    an initial linear time setup, rather than having to recompute the entire tail convergent 
    from scratch for each new convergent.
\end{remark}

\begin{catenadef}[\texttt{SimpleContinuedFraction.tail\_convergent}]
\begin{lstlisting}[style=python]
    from catena import SimpleContinuedFraction

    # All SCFs have a dedicated cache for tail convergents.
    phi = SimpleContinuedFraction(lambda n: 1, integer_part=1)

    phi.cache_handler # CacheHandler for tail convergents; owns the cache and handles caching logic.
    phi.tail_convergent(0) # computes and caches the 0-th tail convergent (1/1)
    phi.tail_convergent(1) # computes and caches the 1-st tail convergent (1/2), reusing the 0-th tail convergent from cache
    
    tail = phi.tail() # the tail is also phi, so it shares the same cache of tail convergents
    tail.cache_handler is phi.cache_handler # True
    tail.tail_convergent(1) # retrieves the 1-st tail convergent from cache (1/2)
\end{lstlisting}
\end{catenadef}

In the literature, tail convergents are typically not introduced explicitly, 
as they are not of primary theoretical interest. However, they provide a useful 
computational tool. In particular, they allow distinct SCFs to share both the 
same generator and the same cache of previously computed values, thereby 
avoiding redundant computations and improving efficiency.

Another point of divergence from the standard presentation is the indexing of 
convergents. In the literature, the $n$-th convergent is often defined as 
$[a_0;\, a_1, \ldots, a_n]$, which corresponds to our $(n-1)$-th convergent. 
We instead adopt zero-based indexing to reflect the conventions of Python, so 
that the first tail convergent corresponds to 
\[
    [0; g(0)] =  \frac{1}{g(0)} = \frac{1}{a_1}.
\]
While minor, this distinction is worth keeping in mind when comparing results 
with the classical literature.

\begin{remark}[RecursionError for Negative Indices]
    The recurrence relations for tail convergents are only defined for $n \geq 0$, and the seed values are defined for $n = -1$ and $n = -2$. 
    Therefore, if the recurrence is called with $n < -2$, it indicates a logic error in the implementation, as it would require access to 
    tail convergents that are not defined. In such cases, a \texttt{RecursionError} should be raised to signal that a negative value has been reached at runtime.
\end{remark}

\subsubsection{Some useful properties of convergents}

\begin{theorem}[Determinant Identity for Convergents\footnote{See, for instance, \cite{khinchin1964}, Theorem 1, p. 4.}]
    Let $x = (a_0, g)$ be a simple continued fraction, where $g: I \to \Zp$. 
    For $n \in I$, let $\frac{p_n}{q_n}$ be the $n$-th convergent of $x$. Then the following identity holds:
    \[
        p_n q_{n-1} - p_{n-1} q_n = (-1)^{n+1} \qquad (n \in I).
    \]
\end{theorem}

\begin{corollary}[Coprimality of Convergents]
    Let $x = (a_0, g)$ be a simple continued fraction, where $g: I \to \Zp$. 
    For $n \in I$, let $\frac{p_n}{q_n}$ be the $n$-th convergent of $x$. Then $p_n$ and $q_n$ are 
    coprime for all $n \in I$; that is,
    \[
        \gcd(p_n, q_n) = 1 \qquad (n \in I).
    \]
\end{corollary}

\begin{proof}
    From the determinant identity,
    \[
        p_n q_{n-1} - p_{n-1} q_n = (-1)^{n+1},
    \]
    any common divisor of $p_n$ and $q_n$ must divide the left-hand side, and hence divide $(-1)^{n+1}$. 
    Therefore, $\gcd(p_n, q_n) = 1$.
\end{proof}

\begin{remark}
    In particular, each convergent is already expressed in lowest terms. This justifies 
    the design choice of not delegating convergent arithmetic to the \texttt{Fraction} 
    class from the Python standard library, which would redundantly attempt to reduce 
    fractions.

    Since the recurrence relations preserve coprimality automatically, it is more 
    efficient to maintain numerators and denominators directly, avoiding both the 
    overhead of normalization and the cost of repeatedly instantiating \texttt{Fraction} 
    objects. Many other algorithms in \catena\ inherit this property as well, allowing 
    coprimality to be preserved without additional (and redundant) computation.
\end{remark}

\begin{corollary}[Difference of Consecutive Convergents]
    \label{coro:difference_consecutive_convergents}
    Let $x = (a_0, g)$ be a simple continued fraction, where $g: I \to \Zp$. 
    For $n \in I$, let $\frac{p_n}{q_n}$ be the $n$-th convergent of $x$. Then the following identity holds:
    \[
        \frac{p_n}{q_n} - \frac{p_{n-1}}{q_{n-1}} = \frac{(-1)^{n+1}}{q_n q_{n-1}} \qquad (n \in I).
    \]
\end{corollary}

\begin{theorem}
    Let $x = (a_0, g)$ be a simple continued fraction, where $g: I \to \Zp$. 
    For $n \in I$, let $\frac{p_n}{q_n}$ be the $n$-th convergent of $x$. Then the following identity holds:
    \[
        p_n q_{n-2} - p_{n-2} q_n = (-1)^n g(n) \qquad (n \in I, n \geq 1).
    \]
\end{theorem}

\begin{corollary}
    Let $x = (a_0, g)$ be a simple continued fraction, where $g: I \to \Zp$. 
    For $n \in I$, let $\frac{p_n}{q_n}$ be the $n$-th convergent of $x$. Then the following identity holds:
    \[
        \frac{p_n}{q_n} - \frac{p_{n-2}}{q_{n-2}} = \frac{(-1)^n g(n)}{q_n q_{n-2}} \qquad (n \in I, n \geq 1).
    \]
\end{corollary}

\begin{theorem}[Monotonicity and Convergence of Convergents]
    Let $x = (a_0, g)$ be a simple continued fraction, where $g: I \to \Zp$. 
    For $n \in I$, let $\frac{p_n}{q_n}$ be the $n$-th convergent of $x$. Then the following properties hold:
    \begin{enumerate}
        \item The sequence of convergents alternates around $x$: for even $n$, we have $c_n > x$, while for odd $n$, we have $c_n < x$.\footnote{
            Due to our indexing convention (where $c_0 = [a_0;\, a_1]$), the parity of the alternating
            inequalities is reversed relative to some standard references.
        }   
        \item The sequence of convergents is strictly decreasing for even indices and strictly increasing for odd indices.
        \item The sequence of convergents converges to $x$ as $n \to \infty$ (if $I = \N$).
    \end{enumerate}
\end{theorem}

\begin{remark}
    The monotonicity and convergence properties of convergents are fundamental to their role as approximations of the original SCF. 
    The alternating nature of the sequence ensures that each new convergent provides a better approximation than the previous one, 
    while the convergence guarantees that we can get arbitrarily close to $x$ by taking sufficiently many convergents.

    This is, among other things, the theoretical basis for being able to define a generic \texttt{\_\_float\_\_} method for SCFs, 
    which is implemented in \catena\ as the value of a fixed convergent call. Namely, \texttt{.convergent(50)} is used as the 
    default approximation for the float value of an arbitrary SCF. The \texttt{machine\_epsilon} constant from the 
    Python standard library is usually around $2^{-52} \approx 2.22 \times 10^{-16}$, so the 50-th 
    convergent of any SCF will be accurate to well within machine precision. A heuristic train of thought for this choice 
    is as follows: since the sequence of convergents alternates around $x$ and converges to $x$, we 
    focus on the value of the difference between the $n$-th convergent and $(n-1)$-th convergent, which, per Corollary \ref{coro:difference_consecutive_convergents}, is given by
    \[
        \left| \frac{p_n}{q_n} - \frac{p_{n-1}}{q_{n-1}} \right| = \frac{1}{q_n q_{n-1}}.
    \]
    On the other hand, the slowest growth rate for the denominators of convergents is given by the Fibonacci sequence, 
    which corresponds to the case where $g(n) = 1$ for all $n$. In this case, we can compute that
    \[
        \frac{1}{q_{50} q_{49}} = \frac{1}{20365011074 \cdot 32951280099} = \frac{1}{671053184118610816326} \approx 1.49 \times 10^{-21}.
    \]
    Therefore, even in the worst-case scenario, the 50-th convergent will be accurate to well within machine precision.
\end{remark}

\begin{catenadef}[\texttt{SimpleContinuedFraction.\_\_float\_\_}]
\begin{lstlisting}[style=python]
from catena import SimpleContinuedFraction

# Example: for phi = [1; 1, 1, 1, ...], the float value is approximately 1.618033988749895
phi = SimpleContinuedFraction(lambda n: 1, integer_part=1)

# We can call float() on phi directly.
float(phi) # approximately 1.618033988749895
\end{lstlisting}
\end{catenadef}

\begin{theorem}[Evaluation from Remainder\footnote{See \cite{khinchin1964}, Theorem 5, p. 7 \& Theorem 7, p. 8.}]
    Let $x = (a_0, g)$ be a simple continued fraction, where $g: I \to \Zp$. 
    For $n \in I$, let $\frac{p_n}{q_n}$ be the $n$-th convergent of $x$, and let $r_n$ be the $n$-th
    remainder of $x$. Then the following identity holds:
    \[
        x = \frac{p_{n-1}r_n + p_{n-2}}{q_{n-1}r_n + q_{n-2}} \qquad (n \in I).
    \]
\end{theorem}

% ─────────────────────────────────────────────────────────────────────────────
\section{Finite Simple Continued Fractions}
% ─────────────────────────────────────────────────────────────────────────────

We now turn to the case of finite simple continued fractions, which correspond to rational numbers.
As mentioned in the definition of the \texttt{SimpleContinuedFraction} class, the body generator 
$g$ is allowed to have a finite domain, in which case the resulting SCF is finite.

Every rational $\nicefrac{p}{q}$ number admits exactly two finite SCF representations, one of which has a final partial quotient of 1.
\[
    \frac{p}{q} = [a_0; a_1, \ldots, a_n] = [a_0; a_1, \ldots, a_{n-1}, a_n - 1, 1] \quad \text{with } a_n \geq 2.
\]

When it is said that a rational number has a \textit{unique} SCF representation, it is understood
that the representation is taken in canonical form, i.e. with $a_n \geq 2$. While this condition
is not enforced at runtime, all constructors that produce finite SCFs (e.g. 
\texttt{FiniteContinuedFraction.from\_rational}) return the canonical form.

\subsection{The Euclidean Algorithm as SCF Expansion}

The SCF coefficients of $p/q$ are precisely the successive quotients produced by the
Euclidean algorithm.
Concretely, \texttt{Convert.from\_rational\_to\_scf} implements:
\begin{align*}
  p_0 &= p, \quad p_1 = q, \\
  a_k &= \left\lfloor \frac{p_k}{p_{k+1}} \right\rfloor, \\
  p_{k+2} &= p_k \bmod p_{k+1}, \qquad k \geq 0,
\end{align*}
until a zero remainder is reached.

\begin{remark}
This correspondence between the Euclidean algorithm and SCF expansion provides
the bijection between finite (canonical) SCFs and rational numbers.

The library uses \texttt{divmod} (Python's built-in, backed by hardware division) for each step,
relying on Python's arbitrary-precision integers so that very large rationals are handled correctly
without overflow.
\end{remark}

\begin{catenadef}[Conversion from Rational-like Types to Finite SCF]
\begin{lstlisting}[style=python]
from catena import FiniteSimpleContinuedFraction

# Convert a rational (Fraction | tuple[int, int] | int) number to its finite SCF representation.
rational_number = (17, 15) # Represents 17/15
finite_scf = FiniteSimpleContinuedFraction.from_rational(rational_number)
print(finite_scf) # FiniteSimpleContinuedFraction(partial_quotients=(7, 2), integer_part=1)

# Convert a float number to its finite SCF representation.
# Note: this conversion may not be exact due to the nature of floating-point representation.
float_number = 3.14
finite_scf_from_float = FiniteSimpleContinuedFraction.from_float(float_number)
print(finite_scf_from_float) # FiniteSimpleContinuedFraction(partial_quotients=(7, 7, 3216856876693, 14), integer_part=3)

finite_scf_from_float_limited = FiniteSimpleContinuedFraction.from_float(3.14, max_denominator=10000) # Limit the denominator to 10,000 to elude floating-point precision artifacts.
print(finite_scf_from_float_limited) # FiniteSimpleContinuedFraction(partial_quotients=(7, 7), integer_part=3)

# Convert a decimal string to its finite SCF representation.
decimal_string = "3.14"
finite_scf_from_decimal = FiniteSimpleContinuedFraction.from_decimal(decimal_string)
print(finite_scf_from_decimal) # FiniteSimpleContinuedFraction(partial_quotients=(7, 7), integer_part=3)
\end{lstlisting}
\end{catenadef}

\begin{remark}
  \texttt{FiniteSimpleContinuedFraction.from\_decimal} takes a decimal string as input rather than a 
  \texttt{decimal.Decimal} object. This is because the latter depends on a global precision context 
  (\texttt{getcontext().prec}), and operations such as \texttt{Decimal(2).sqrt()} produce approximations 
  whose accuracy depends on that context. As a result, converting directly from \texttt{Decimal} would 
  require an additional precision parameter and would not be guaranteed to produce consistent results 
  across different environments. Using a decimal string avoids this ambiguity by making the input 
  representation explicit in the API call.
\end{remark}

% ─────────────────────────────────────────────────────────────────────────────
\section{Convergents and the Three-Term Recurrence}
% ─────────────────────────────────────────────────────────────────────────────

\subsection{Definition}

\begin{definition}[Convergent]
  The $k$-th \emph{convergent} of $[a_0; a_1, \ldots, a_n]$ is
  $x_k = p_k / q_k$ where the numerator--denominator pairs satisfy the recurrence
  \begin{equation}\label{eq:recurrence}
    \begin{pmatrix} p_k \\ q_k \end{pmatrix}
    = a_k \begin{pmatrix} p_{k-1} \\ q_{k-1} \end{pmatrix}
    + \begin{pmatrix} p_{k-2} \\ q_{k-2} \end{pmatrix}
  \end{equation}
  with initial conditions $p_{-1} = 1$, $p_0 = a_0$, $q_{-1} = 0$, $q_0 = 1$.
\end{definition}

\subsection{Implementation}

\texttt{FinateSimpleContinuedFraction.\_convergent} implements \cref{eq:recurrence}
recursively, re-indexed so that \texttt{body[0]} $= a_1$, \texttt{body[k]} $= a_{k+1}$.
The base cases in the code are:
\begin{align*}
  n=0: &\quad (p_0, q_0) = (1,\; a_1),  \\
  n=1: &\quad (p_1, q_1) = (a_2,\; a_1 a_2 + 1).
\end{align*}

\begin{remark}[Design decision: body-indexed convergents]
The \texttt{convergent(n)} method indexes into \texttt{body}, so $n$ ranges from $0$ to
$\texttt{size}-1$ where $\texttt{size} = |\texttt{body}|$.
The actual rational value requires adding \texttt{head} to $p/q$, which is what
\texttt{convergent\_as\_float} does.
This separation keeps the convergent arithmetic clean and free of a sign-carrying integer part.
\end{remark}

\subsection{Key Properties}

\begin{proposition}[Determinant formula]
  For any $k \geq 0$:
  \[
    p_{k-1} q_k - p_k q_{k-1} = (-1)^k.
  \]
  In particular $\gcd(p_k, q_k) = 1$ for all $k$.
\end{proposition}

\begin{proposition}[Best rational approximation]
  The convergents $p_k/q_k$ are the \emph{best rational approximations} to the
  represented rational in the sense that no fraction with denominator $\leq q_k$
  approximates the value better than $p_k/q_k$.
\end{proposition}

\begin{remark}
  The best-approximation property motivates using SCF convergents for high-precision
  string output: \texttt{convergent\_as\_string} and \texttt{convergent\_float\_as\_string}
  use \texttt{StringMethod.string\_divition}, which computes a long-division decimal
  expansion of a convergent pair $(p_k, q_k)$ to an arbitrary number of digits.
\end{remark}

% ─────────────────────────────────────────────────────────────────────────────
\section{Caching Strategy}
% ─────────────────────────────────────────────────────────────────────────────

\subsection{Motivation}

Two separate caching concerns arise in \catena.

\textbf{Tail-convergent caching.}  The recurrence
\[
    h_n = g(n)\,h_{n-1} + h_{n-2}, \quad k_n = g(n)\,k_{n-1} + k_{n-2}
\]
has optimal substructure: $h_n$ depends on $h_{n-1}$ and $h_{n-2}$.
Without memoisation a naive recursive implementation recomputes all earlier
tail convergents each time a later one is requested.  With memoisation each
value is computed once, reducing the total cost from $O(2^n)$ to $O(n)$.
All \texttt{SimpleContinuedFraction} instances carry a dedicated
\texttt{\_tail\_cache} for this purpose.

\textbf{Generator caching.}  When the partial-quotient function $g: \N \to \Zp$
is expensive to evaluate (e.g.\ a number-theoretic function or a transcendental
constant computed digit by digit), re-evaluating $g(n)$ on every recurrence step
would dominate the runtime.  \texttt{CachedGenerator} wraps an arbitrary callable
and memoises each distinct index in a \texttt{BaseCache}.

\subsection{The \texttt{BaseCache} Primitive}

\catena's caching layer is built around a single primitive defined in
\texttt{catena.cache}: \texttt{BaseCache}.  It is a self-computing,
append-only dictionary (\texttt{UserDict} subclass) that unifies the three
separate classes that existed in earlier versions (\texttt{Cache},
\texttt{OrdinalCache}, and \texttt{CacheHandler}).

\begin{description}[style=nextline]
  \item[Immutability after write.]
    Once a key is stored, any subsequent attempt to overwrite it raises
    \texttt{KeyError}.  This is correct for convergents and generator values
    alike, because both are deterministic functions of their index for a fixed SCF.

  \item[Self-computing on miss.]
    \texttt{BaseCache} is constructed with a bound callable \texttt{func}.
    Accessing a missing key via \texttt{cache[key]} calls \texttt{func(key)},
    stores the result, and returns it — no external decorator or wrapper is
    needed.  Explicit writes (\texttt{cache[key] = value}) are also accepted
    and bypass the function call.

  \item[Key-range tracking.]
    For integer keys, \texttt{smallest\_key} and \texttt{largest\_key}
    are maintained without iterating the dictionary.
    The \texttt{largest\_key} sentinel drives the tail-convergent iterative
    fill: the recurrence advances from \texttt{largest\_key + 1} to $n$,
    computing only the truly missing entries in $O(\text{gap})$ time and
    $O(1)$ call-stack depth regardless of $n$.

  \item[Usage statistics.]
    \texttt{call\_count} counts cache misses (calls to \texttt{func}),
    \texttt{read\_count} counts cache hits (stored-result retrievals), and
    \texttt{write\_count} counts all writes (including those from seeding
    and intermediate fills).

  \item[Capacity management.]
    An optional \texttt{maxsize} parameter caps the number of stored entries.
    When a new write would exceed the limit, \texttt{prune(n, order)} is
    called automatically using a configurable \texttt{prune\_key} callable.
    The default policy retains the \texttt{maxsize\,-\,1} largest-keyed
    entries, keeping the recurrence frontier intact.  \texttt{reset()} clears
    the cache and all counters; \texttt{clear()} clears entries only, preserving
    \texttt{call\_count} and \texttt{read\_count}.

  \item[Construction-time seeding.]
    The \texttt{seed} constructor parameter pre-populates the cache from a
    plain dict without incrementing any statistics counters.  This is the
    mechanism used by \texttt{CachedGenerator.advance} and
    \texttt{CachedGenerator.insert} to carry relevant entries across
    cache-shifting operations.

  \item[Blocked mutation methods.]
    The mutating \texttt{UserDict} methods (\texttt{pop}, \texttt{popitem},
    \texttt{\_\_delitem\_\_}, \texttt{update}, \texttt{setdefault},
    \texttt{fromkeys}, \texttt{\_\_or\_\_}) are all disabled and raise
    \texttt{NotImplementedError}.  Use \texttt{prune} or \texttt{reset} to
    reduce cache size.
\end{description}

\subsection{Generator-Level Caching: \texttt{CachedGenerator}}

\texttt{CachedGenerator} is a subclass of \texttt{Generator} that internally
holds a \texttt{BaseCache} bound to the wrapped callable.  Every call
\texttt{cgen(n)} is routed through the cache: on a miss the underlying function
is invoked and the result stored; on a hit the stored result is returned directly.

\textbf{When to use each generator type.}
\begin{enumerate}
  \item \texttt{Generator}: the default.  Wrap any callable $g: \N \to \Zp$
        that is cheap to evaluate (e.g.\ a closed-form formula).
  \item \texttt{CachedGenerator}: use when $g(n)$ is expensive \emph{and}
        catena's cache introspection is needed (\texttt{call\_count},
        \texttt{read\_count}, \texttt{stats}, \texttt{BaseCache}
        bookkeeping).  Also the right choice when one expensive generator
        must be shared across several SCF instances so that all instances
        draw from the same memoised pool.
  \item Sharing via \texttt{Generator} wrapping an \texttt{@lru\_cache}-decorated
        function achieves the same functional sharing without catena's
        introspection layer.  Use this pattern for recursive functions
        (Fibonacci, etc.): apply \texttt{@functools.lru\_cache} to the function
        itself \emph{before} passing it to \texttt{Generator} or
        \texttt{CachedGenerator}, because \texttt{CachedGenerator} only
        intercepts top-level calls and does not propagate into self-recursive
        calls inside the function body.
  \item \texttt{FiniteGenerator}: use for fixed finite sequences; no caching
        is needed because lookup is $O(1)$ by array index.
\end{enumerate}

\begin{catenadef}[Caching Patterns]
\begin{lstlisting}[style=python]
from catena.generators import Generator, CachedGenerator
from catena import SimpleContinuedFraction
import functools, math

# --- Pattern 1: expensive non-recursive generator ---
# g(n) = number of divisors of (n+2) not exceeding sqrt(n+2)
def small_divisors(n):
    m = n + 2
    return sum(1 for d in range(1, math.isqrt(m) + 1) if m % d == 0)

cgen = CachedGenerator(small_divisors)
scf = SimpleContinuedFraction(cgen, integer_part=0)
scf.convergent(20)       # fills cache entries 0..20 as a side-effect
cgen.call_count          # 21  (one miss per index)
cgen.read_count          # 0   (no hits yet from this SCF)

# --- Pattern 2: sharing one cache across multiple SCFs ---
scf2 = SimpleContinuedFraction(cgen, integer_part=1)
scf2.convergent(10)      # cgen(0)..cgen(10) are already cached
cgen.read_count          # 11  (all hits: entries 0..10 were populated by scf)

# --- Pattern 3: recursive generator - pre-memo with lru_cache ---
@functools.lru_cache(maxsize=None)
def fib(n):
    return fib(n - 1) + fib(n - 2) if n > 1 else n + 1

# Wrap in CachedGenerator to add BaseCache introspection on top
cgen_fib = CachedGenerator(fib)     # fib's own lru_cache handles recursion
cgen_fib(10)                        # fib(10) was already cached by lru_cache
cgen_fib.call_count                 # 1
cgen_fib.cache.reset()              # clears BaseCache; lru_cache unchanged
\end{lstlisting}
\end{catenadef}

\subsection{Shared Caches and the Tail}

All \texttt{tail()}, \texttt{shift()}, and \texttt{segment(size)} operations
create a new SCF instance via \texttt{\_from\_shared}, which copies the
\emph{reference} to \texttt{\_tail\_cache} rather than duplicating the cache
object.  Consequently both the original and the derived SCF read from and write
to the same \texttt{BaseCache}; a tail-convergent computed by one is
immediately available to the other at zero cost.

This is the mechanism that makes the following hold:

\begin{catenadef}[Shared tail-convergent cache]
\begin{lstlisting}[style=python]
from catena import SimpleContinuedFraction

phi = SimpleContinuedFraction(lambda n: 1, integer_part=1)
phi.tail_convergent(9)             # populates cache entries 0..9

tail = phi.tail()                  # new instance, same _tail_cache
tail.cache_handler is phi.cache_handler   # True
tail.tail_convergent(9)            # O(1) cache hit -- no recomputation
phi.cache_handler.read_count       # incremented by the hit above
\end{lstlisting}
\end{catenadef}



% ─────────────────────────────────────────────────────────────────────────────
\section{The Greedy (Sylvester) Algorithm}
% ─────────────────────────────────────────────────────────────────────────────

\texttt{Method.greedy\_algorithm} decomposes a fraction $p/q$ into an
\emph{Egyptian fraction} — a sum of distinct unit fractions — using the greedy rule
\[
  \frac{p}{q} = \frac{1}{\lceil q/p \rceil} + \frac{p\lceil q/p \rceil - q}{q\lceil q/p \rceil}.
\]
This is distinct from the SCF expansion; it is sometimes called the Sylvester sequence.
The intermediate use of \texttt{Method.simplify} at each step ensures the process terminates.

\begin{remark}[Relevance to SCF]
  The greedy algorithm appears in \texttt{catena} via \texttt{Convert.from\_float\_to\_armonic}
  and \texttt{Convert.from\_str\_to\_armonic}, providing an alternative representation of a
  rational that may be compared to the SCF expansion.
  Their relationship to the Stern--Brocot tree and the Calkin--Wilf tree is a topic for a
  future section.
\end{remark}

% ─────────────────────────────────────────────────────────────────────────────
\section{Planned Extensions}
% ─────────────────────────────────────────────────────────────────────────────

The following are flagged as future development areas.
Each will warrant its own section in this document.

\subsection{Periodic Continued Fractions (Quadratic Irrationals)}

By Lagrange's theorem, $\alpha \in \R$ has a periodic SCF if and only if
$\alpha$ is a quadratic irrational.
Planned class: \texttt{PeriodicSimpleContinuedFraction}.

\subsection{Arithmetic on Continued Fractions}

Gosper (1972) showed how to perform the four arithmetic operations on SCFs
in a streaming fashion using $2 \times 2$ and $2 \times 2 \times 2$ integer matrices,
without fully evaluating either operand.
This is the natural future direction for \texttt{FinateSimpleContinuedFraction} arithmetic.

\cite{gosper1972} % placeholder until refs.bib is populated

\subsection{Generalised Continued Fractions}

The general form $b_0 + K_{n=1}^{\infty} a_n / b_n$ with arbitrary integer or
function sequences.
Applications: rapid evaluation of $\pi$, $e$, $\ln 2$, Bessel functions, etc.

\subsection{Metric Theory and Gauss Map}

The map $T: x \mapsto \{1/x\}$ (fractional part of $1/x$) on $(0,1)$ is the
\emph{Gauss map}; its invariant measure $\mu(A) = \frac{1}{\ln 2}\int_A \frac{dx}{1+x}$
governs the distribution of partial quotients (Gauss--Kuzmin statistics).
This will be relevant for statistical analysis routines.

\subsection{Multidimensional Generalisations}

Jacobi--Perron and related algorithms for simultaneous approximation; connection to
lattice reduction (LLL algorithm).

% ─────────────────────────────────────────────────────────────────────────────
\section{Notation Index}
% ─────────────────────────────────────────────────────────────────────────────

\begin{tabular}{ll}
\toprule
Symbol & Meaning \\
\midrule
$[a_0; a_1, \ldots, a_n]$ & Finite simple continued fraction \\
$(p_k, q_k)$               & $k$-th convergent numerator/denominator \\
\texttt{head}              & $a_0$ (integer part) \\
\texttt{body}              & $(a_1, \ldots, a_n)$ (partial quotients) \\
\texttt{size}              & $n = |\texttt{body}|$ \\
$T$                        & Gauss map $x \mapsto \{1/x\}$ \\
\bottomrule
\end{tabular}


% ─────────────────────────────────────────────────────────────────────────────
\printbibliography
% ─────────────────────────────────────────────────────────────────────────────

\end{document}
